\chapter{Krajobraz regulacyjny}

\section{Krajowe ramy prawne cyberbezpieczeństwa i~ochrony danych}

Polska, jako państwo członkowskie Unii Europejskiej, od lat aktywnie uczestniczy w budowaniu europejskiego systemu ochrony danych i~cyberbezpieczeństwa. Punktem wyjścia dla krajowych regulacji w~tym zakresie stało się wdrożenie unijnej dyrektywy NIS z~2016 roku \autocite{dyrektywa_nis}, która zobowiązała państwa członkowskie do ustanowienia krajowych norm cyberbezpieczeństwa. W~odpowiedzi na te wymagania, Polska uchwaliła ustawę o~krajowym systemie cyberbezpieczeństwa, przyjętą przez Sejm RP dnia 5~lipca 2018~roku i~podpisaną przez Prezydenta RP dnia 1~sierpnia 2018~r., która weszła w~życie 28~sierpnia 2018~r.
\autocite{uksc2018}.

\subsection{Zakres obowiązków wynikający z ustawy o~krajowym systemie cyberbezpieczeństwa z~2018 r.}

Ustawa o~krajowym systemie cyberbezpieczeństwa~\autocite{uksc2018} nakłada kompleksowe obowiązki przede wszystkim na operatorów usług kluczowych (\textit{OUK}) oraz dostawców usług cyfrowych (\textit{DUC}), działających w~sektorach takich jak energetyka, transport, bankowość, ochrona zdrowia czy infrastruktura cyfrowa.

\begin{table}[H]
\centering
\caption{Główne obowiązki nałożone ustawą o~KSC z~2018~r.}
\label{tab:ksc_obowiazki}
\begin{tabular}{>{\bfseries}p{4cm} p{9.5cm}}
\hline
Obszar obowiązku & Opis \\
\hline
Zarządzanie ryzykiem
  & Regularna identyfikacja i~ocena ryzyk cyberbezpieczeństwa; opracowanie
    i~wdrożenie strategii ochrony adekwatnej do specyfiki działalności. \\[4pt]
Środki techniczne i~org.
  & Zapewnienie bezpiecznej eksploatacji systemów, w~tym bezpieczeństwa
    fizycznego, kontroli dostępu oraz ciągłości świadczenia usług. \\[4pt]
Monitorowanie i~zgłaszanie incydentów
  & Bieżące monitorowanie systemów informatycznych; raportowanie
    incydentów do właściwych zesp.\ CSIRT (Zespół reagowania na incydenty bezpieczeństwa komputerowego)
    (CSIRT GOV, CSIRT NASK, CSIRT MON). \\[4pt]
Dokumentacja
  & Opracowanie i~przechowywanie przez co~najmniej 2~lata dokumentacji
    dotyczącej cyberbezpieczeństwa systemu informacyjnego. \\[4pt]
Audyty
  & Przeprowadzanie regularnych audytów bezpieczeństwa oraz szkoleń
    dla pracowników. \\[4pt]
Współpraca z~organami
  & Koordynacja działań z~organami właściwymi ds.\ cyberbezpieczeństwa;
    udostępnianie wyników audytów w~razie kontroli. \\
\hline
\end{tabular}
\end{table}

\subsection{Nowelizacja UKSC z~23 stycznia 2026~roku\ (implementacja NIS2)}

W~dniu 23~stycznia 2026~r.\ Sejm uchwalił nowelizację ustawy o~krajowym systemie cyberbezpieczeństwa~\autocite{uksc2026}, stanowiącą implementację
unijnej dyrektywy NIS2~\autocite{dyrektywa_nis2} do polskiego porządku prawnego. Ustawa została następnie podpisana przez Prezydenta RP, który jednocześnie
skierował ją do Trybunału Konstytucyjnego, podnosząc wątpliwości dotyczące m.in.\ obowiązku wymiany sprzętu bez odszkodowania oraz sposobu podejmowania
decyzji przez organy nadzorcze.

\begin{table}[H]
\centering
\caption{Kluczowe zmiany wprowadzone nowelizacją UKSC z~2026~r.}
\label{tab:ksc_nowelizacja}
\begin{tabular}{>{\bfseries}p{5cm} p{8.5cm}}
\hline
Zmiana & Opis \\
\hline
Rozszerzenie zakresu podmiotowego
  & Poszerzenie katalogu branż i~obniżenie progu wejścia do
    systemu \\[4pt]
Obowiązek samoidentyfikacji
  & Organizacja samodzielnie  ma obowiązek ustalić swój
    status, dokonać kwalifikacji i~złożyć wniosek o~wpis do
    właściwego wykazu. \\[4pt]
Odpowiedzialność zarządcza
  & Odpowiedzialność przestaje być wyłącznie operacyjna i~staje się
    elementem ładu korporacyjnego oraz osobistej odpowiedzialności
    członków zarządu. \\[4pt]
Kary finansowe
  & Podmioty kluczowe: do 10~mln~EUR lub 2\% rocznego światowego obrotu;
    podmioty ważne: do 7~mln~EUR lub 1,4\% obrotu. \\[4pt]
Odroczenie kar
  & Kary pieniężne za naruszenia mogą być nakładane dopiero po upływie
    2~lat od wejścia w~życie ustawy. \\[4pt]
Modyfikacje definicji
  & Zawężenie definicji organizacji badawczej i~doprecyzowanie zakresu
    podmiotowego wobec środowisk naukowych. \\
\hline
\end{tabular}
\end{table}

\subsection{Ustawa o~ochronie danych osobowych z~10 maja 2018~r.}

Z punktu widzenia bezpieczeństwa, istotną ustawą w~krajowym prawodawstwie jest Ustawa z~dnia 10~maja 2018~r.\
o~ochronie danych osobowych~\autocite{uodo2018}, która stanowi akt wykonawczy wobec rozporządzenia RODO~\autocite{rodo} i~reguluje m.in.\ zasady działania Prezesa Urzędu Ochrony Danych Osobowych (PUODO) jako krajowego organu nadzorczego, tryb postępowania przed tym organem, szczegółowe zasady przetwarzania danych oraz krajowe przepisy uzupełniające RODO tam, gdzie rozporządzenie
pozostawia swobodę regulacyjną państwom członkowskim.

\subsubsection{Odpowiedzialność administracyjna na gruncie RODO}

Rozporządzenie RODO~\autocite{rodo} przewiduje dwustopniowy system administracyjnych kar
pieniężnych (art.~83 RODO).

\begin{table}[H]
\centering
\caption{Administracyjne kary pieniężne na mocy art.~83 RODO}
\label{tab:rodo_kary}
\begin{tabular}{>{\bfseries}p{4.5cm} p{9cm}}
\hline
Wysokość kary & Przykładowe naruszenia \\
\hline
Do 10~mln~EUR lub 2\% rocznego obrotu
  & Brak rejestru czynności przetwarzania; brak wyznaczenia inspektora ochrony
    danych w~wymaganych przypadkach; brak współpracy z~organem nadzorczym. \\[4pt]
Do 20~mln~EUR lub 4\% rocznego obrotu
  & Złamanie podstawowych zasad przetwarzania danych; naruszenie praw osób,
    których dane dotyczą; nielegalny transfer danych do państw trzecich. \\
\hline
\end{tabular}
\end{table}

Organ nadzorczy dysponuje również wachlarzem środków nacisku (art.~58 RODO): ostrzeżenia i~upomnienia, nakazanie spełnienia żądań osoby,
której dane dotyczą, nakazanie dostosowania operacji przetwarzania do przepisów RODO, a~w~najpoważniejszych przypadkach, czasowe lub całkowite ograniczenie
przetwarzania, w~tym zakaz przetwarzania danych oraz nakaz ich usunięcia lub powiadomienia o~naruszeniu.

\subsection{Ustawa o~zwalczaniu nieuczciwej konkurencji (1993)}

Ustawa z~dnia 16~kwietnia 1993~r.\ o~zwalczaniu nieuczciwej konkurencji~\autocite{uznk1993} stanowi istotne uzupełnienie systemu ochrony danych, zwłaszcza w~zakresie ochrony tajemnicy przedsiębiorstwa. Zgodnie z~art.~11 u.z.n.k., czynem nieuczciwej konkurencji jest ujawnienie, wykorzystanie lub pozyskanie cudzych informacji stanowiących tajemnicę przedsiębiorstwa. Za tajemnicę przedsiębiorstwa uznaje się informacje techniczne,
technologiczne, organizacyjne lub inne posiadające wartość gospodarczą, które nie są powszechnie znane i~co do których przedsiębiorca podjął należyte środki
w~celu zachowania ich poufności.

W~kontekście naruszeń danych i~cyberbezpieczeństwa ustawa ta nabiera szczególnego znaczenia: wyciek danych stanowiących tajemnicę przedsiębiorstwa (np.\ w~wynik ataku hakerskiego, działania nielojalnego pracownika lub wadliwego zabezpieczenia systemów) może jednocześnie stanowić naruszenie RODO oraz czyn nieuczciwej konkurencji z~u.z.n.k. Ustawa przewiduje cywilnoprawne środki ochrony, w~tym roszczenia o~zaniechanie, usunięcie skutków, naprawienie szkody i~wydanie bezpodstawnie uzyskanych korzyści~\autocite{uznk1993}.

\subsection{Odpowiedzialność cywilna wg kodeksu cywilnego}

Odpowiedzialności z~tytułu naruszenia przepisów o~ochronie danych nie można rozpatrywać wyłącznie przez pryzmat RODO, kluczową rolę odgrywa tu również
Kodeks cywilny~\autocite{kc1964}.

\begin{table}[H]
\centering
\caption{Reżimy odpowiedzialności cywilnej w~obszarze ochrony danych}
\label{tab:odpowiedzialnosc_cywilna}
\begin{tabular}{>{\bfseries}p{5cm} p{8.5cm}}
\hline
Podstawa prawna & Opis \\
\hline
Art.~415 k.c.: odpowiedzialność deliktowa
  & Administrator lub podmiot przetwarzający, który ze swojej winy wyrządził
    szkodę osobie fizycznej poprzez naruszenie zasad ochrony danych,
    zobowiązany jest do jej naprawienia. Przesłankami są: bezprawne działanie
    lub zaniechanie, wina, szkoda oraz związek przyczynowo-skutkowy. \\[4pt]
Art.~23, 24 i~448 k.c.: dobra osobiste
  & Prywatność jest traktowana jako dobro osobiste. W~razie jego naruszenia
    pokrzywdzony może żądać zaniechania naruszeń i~zasądzenia
    zadośćuczynienia (sądy orzekają w~zbiegu z~art.~82 RODO). \\[4pt]
Art.~82 ust.~4 RODO w~zw.\ z~k.c.: solidarność
  & Jeżeli w~tym samym przetwarzaniu uczestniczy więcej niż jeden podmiot,
    odpowiadają oni solidarnie za całą szkodę, co zapewnia poszkodowanemu
    rzeczywiste uzyskanie odszkodowania. \\[4pt]
Zbieg podstaw odpowiedzialności
  & Nałożenie kary przez PUODO nie zwalnia administratora z~równoległej
    odpowiedzialności cywilnej wobec osób poszkodowanych. \\
\hline
\end{tabular}
\end{table}

\subsection{Odpowiedzialność karna: Kodeks karny}

Odrębnym, lecz komplementarnym reżimem jest odpowiedzialność karna przewidziana w~Kodeksie karnym~\autocite{kk1997}, w~szczególności w~rozdziale XXXIII
,,Przestępstwa przeciwko ochronie informacji'' (art.~265--269b k.k.).

\begin{table}[H]
\centering
\caption{Przestępstwa z~zakresu ochrony informacji (art.~265--269b k.k.)}
\label{tab:kk_przestepstwa}
\begin{tabular}{>{\bfseries}p{3.5cm} p{10cm}}
\hline
Przepis & Opis \\
\hline
Art.~267 k.k.\: hacking
  & Kara grzywny, ograniczenia wolności albo pozbawienia wolności do 2~lat
    za bezprawne uzyskanie dostępu do informacji poprzez przełamanie
    zabezpieczeń elektronicznych, podpięcie do sieci lub stosowanie
    oprogramowania szpiegującego; penalizowane również ujawnianie
    tak pozyskanych informacji. \\[4pt]
Art.~268 i~268a k.k. & Niszczenie lub uszkadzanie danych informatycznych i~udaremnianie osobie
    uprawnionej zapoznania się z~informacją; kara do 3~lat pozbawienia
    wolności, a~przy znacznej szkodzie majątkowej od 3~miesięcy do 5~lat. \\[4pt]
Art.~269 k.k.: sabotaż komputerowy
  & Uszkodzenie danych informatycznych lub oprogramowania o~szczególnym
    znaczeniu dla obronności lub bezpieczeństwa państwa;
    kara pozbawienia wolności od 6~miesięcy do 8~lat. \\[4pt]
Art.~269a k.k.: zakłócenie systemu
  & Istotne zakłócenie pracy systemu komputerowego lub sieci
    teleinformatycznej przez transmisję, usunięcie lub zmianę danych;
    kara pozbawienia wolności od 3~miesięcy do 5~lat. \\[4pt]
Art.~269b k.k.: narzędzia ataków
  & Wytwarzanie, pozyskiwanie lub udostępnianie narzędzi informatycznych
    przeznaczonych do popełnienia powyższych przestępstw. \\[4pt]
Art.~23 u.z.n.k.\ (uzupełniająco)
  & Ujawnienie tajemnicy przedsiębiorstwa przez osobę zobowiązaną do jej
    zachowania, jeżeli wyrządza to poważną szkodę: kara grzywny, ograniczenia
    wolności albo pozbawienia wolności do 2~lat. \\
\hline
\end{tabular}
\end{table}

Warto podkreślić, że reżim karny funkcjonuje równolegle do odpowiedzialności administracyjnej (kary RODO/UKSC) i~cywilnej, skazanie sprawcy nie eliminuje
roszczeń odszkodowawczych poszkodowanych ani nie uchyla administracyjnych kar nałożonych na administratora danych~\autocite{kk1997}\autocite{rodo}.


\subsection{Dyrektywa NIS2 jako fundament europejskiego cyberbezpieczeństwa}

Dyrektywa Parlamentu Europejskiego i~Rady (UE) 2022/2555 z dnia 14~grudnia 2022~r., znana jako NIS2~\autocite{dyrektywa_nis2}, reformuje unijny system cyberbezpieczeństwa. NIS2 uchyliła i~zastąpiła dyrektywę NIS z~2016~r.~\autocite{dyrektywa_nis}, która  okazała się niewystarczająca wobec  wzrostu zagrożeń cybernetycznych oraz rosnącej zależności gospodarki i~administracji od systemów informatycznych.
Dyrektywa NIS2 weszła w życie 16~stycznia 2023~r., zaś państwa członkowskie zobowiązane były do dokonania jej transpozycji do krajowych porządków prawnych
do 17~października 2024~r. W Polsce implementacji dyrektywy dokonała nowelizacja ustawy o Krajowym Systemie Cyberbezpieczeństwa z~dnia 23~stycznia 2026~r.~\autocite{uksc2026}.

\subsection*{Zakres podmiotowy: podmioty kluczowe i~ważne}

Dyrektywa NIS2 odchodzi od dotychczasowego podziału na operatorów usług kluczowych
i dostawców usług cyfrowych, dyrektywa definiuje nowe kategorie: {podmioty kluczowe} (ang.\ \textit{essential entities}) oraz {podmioty ważne}
(ang.\ \textit{important entities}). 
Dyrektywą objęte są podmioty będące co najmniej średnim przedsiębiorstwem wg definicji MŚP\autocite{deficja_msp_2018} (tj.\ zatrudniające min. 50 osób i~osiągające obrót roczny powyżej 10~mln~EUR), działające w sektorach wskazanych w załącznikach i~i II do dyrektywy. Jednocześnie przewidziano szereg wyjątków,
na mocy których dyrektywa stosuje się niezależnie od wielkości podmiotu.

\begin{table}[H]
\centering
\caption{Sektory objęte dyrektywą NIS2 według kategorii podmiotów}
\label{tab:nis2_sektory}
\begin{tabular}{>{\bfseries}p{4.5cm} p{8.5cm}}
\hline
Kategoria & Sektory (wybrane) \\
\hline
Podmioty kluczowe (załącznik I)
 & Energetyka (energia elektryczna, gaz, ropa, wodór); transport (lotniczy,
    kolejowy, wodny, drogowy); bankowość i~infrastruktura rynków finansowych;
    ochrona zdrowia; woda pitna i~ścieki; infrastruktura cyfrowa (IXP, DNS,
    usługi chmurowe, centra danych, sieci CDN); zarządzanie usługami ICT (B2B);
    administracja publiczna; przestrzeń kosmiczna. \\[4pt]
Podmioty ważne (załącznik II)
 & Usługi pocztowe i~kurierskie; gospodarka odpadami; produkcja
    i~dystrybucja chemikaliów; produkcja i~dystrybucja żywności; produkcja
    (wyroby medyczne, sprzęt elektroniczny, pojazdy); dostawcy usług cyfrowych
    (platformy handlowe, wyszukiwarki, serwisy społecznościowe); organizacje
    badawcze. \\
\hline
\end{tabular}
\end{table}

Zakresem dyrektywy nie są objęte podmioty administracji publicznej prowadzące działalność w dziedzinie bezpieczeństwa narodowego, obronności lub ścigania
przestępstw, podmioty wyłączone ze stosowania rozporządzenia DORA~\autocite{dora2022} (stosownie do art.~2 ust.~4 DORA), a~także dostawcy świadczący usługi
wyłącznie na rzecz zwolnionych organów administracji publicznej. Z obowiązku zwolnienia nie można jednak skorzystać wobec dostawców usług zaufania.

\subsection*{Obowiązki podmiotów kluczowych i~ważnych}

Dyrektywa NIS2 wprowadza identyczne obowiązki dla obu kategorii podmiotów (kluczowe oraz ważne), rozróżniając jedynie intensywność nadzoru i~poziom kar. Obowiązki te mają być stosowane w sposób proporcjonalny do specyfiki, wielkości i~poziomu narażenia danego podmiotu na ryzyko.

Podmioty muszą zaimplementować \textit{System Zarządzania Bezpieczeństwem Informacji}, w tym zarządzania ryzykiem. Wymaga do opracowania, wdrożenia i~testowania planów ciągłości działania (\textit{BCP}) oraz politykę bezpieczeństwa łańcucha dostaw, polegająceł na szacowaniu ryzk współpracy w dostawcami oraz poddostawcami.

Każdy podmiot kluczowy i~ważny zobowiązany jest do wdrożenia odpowiednich i~proporcjonalnych środków technicznych, operacyjnych i~organizacyjnych,
obejmujących co najmniej:

\begin{table}[H]
\centering
\caption{Minimalne środki zarządzania ryzykiem wymagane przez art.~21 NIS2}
\label{tab:nis2_ryzyko}
\begin{tabular}{>{\bfseries}p{4.5cm} p{8.5cm}}
\hline
Obszar & Wymagania minimalne \\
\hline
Polityki bezpieczeństwa i~analizy ryzyka
  & Opracowanie i~wdrożenie polityk bezpieczeństwa systemów informatycznych;
    regularna ocena ryzyka uwzględniająca zagrożenia, prawdopodobieństwo
    i~dotkliwość incydentów oraz skutki społeczne i~gospodarcze. \\[4pt]
Obsługa incydentów
  & Ustanowienie procedur wykrywania, analizy, reakcji i~odtwarzania
    po incydentach bezpieczeństwa. \\[4pt]
Ciągłość działania
  & Wdrożenie planów ciągłości działania (BCP) i~odtwarzania po awarii (DRP),
    w tym zarządzanie kopiami zapasowymi i~zarządzanie kryzysowe. \\[4pt]
Bezpieczeństwo łańcucha dostaw
  & Ocena i~zarządzanie ryzykiem ze strony dostawców i~usługodawców ICT;
    stosowanie klauzul bezpieczeństwa w umowach z dostawcami. \\[4pt]
Bezpieczeństwo sieci i~systemów
  & Środki kontroli dostępu, uwierzytelnianie wieloskładnikowe (MFA),
    szyfrowanie danych w spoczynku i~w tranzycie, segmentacja sieci,
    zarządzanie podatnościami i~aktualizacjami. \\[4pt]
Szkolenia i~świadomość
  & Regularne szkolenia dla zarządu i~pracowników; organy zarządzające
    obowiązane do zatwierdzania środków zarządzania ryzykiem i~odbycia
    wymaganych szkoleń. \\[4pt]
Kryptografia i~szyfrowanie
  & Wdrożenie polityk stosowania kryptografii i~szyfrowania w odniesieniu
    do przetwarzanych danych i~komunikacji. \\[4pt]
Bezpieczeństwo zasobów ludzkich
  & Procedury zarządzania dostępem i~polityki kontroli dostępu oparte
    na zasadzie najmniejszych uprawnień; bezpieczeństwo personelu. \\
\hline
\end{tabular}
\end{table}

Środki zarządzania ryzykiem muszą być zatwierdzone przez organ zarządzający podmiotu. W przypadku stwierdzenia niezgodności z wymaganiami dyrektywy,
podmiot zobowiązany jest do bezzwłocznego wdrożenia środków naprawczych.

\subsubsection{Zgłaszanie incydentów (art.~23 NIS2)}

Podmioty kluczowe i~ważne zobowiązane są do zgłaszania {poważnych incydentów} do właściwego zespołu CSIRT lub organu nadzorczego. Za poważny
incydent uznaje się zdarzenie, które powoduje lub może spowodować znaczne zakłócenia operacyjne lub straty finansowe, albo wpływa na inne osoby lub
organizacje powodując znaczne szkody.

\begin{table}[H]
\centering
\caption{Harmonogram obowiązkowego raportowania incydentów na mocy art.~23 NIS2}
\label{tab:nis2_raportowanie}
\begin{tabular}{>{\bfseries}p{4cm} p{9.5cm}}
\hline
Termin & Zakres informacji \\
\hline
Max. 24 h:  wczesne ostrzeżenie
  & Niezwłoczne poinformowanie o poważnym incydencie z~wskazaniem: czy
    incydent mógł być wywołany działaniem bezprawnym lub złym zamiarem;
    czy mógł wywrzeć wpływ transgraniczny. \\[4pt]
Max. 72 h:  zgłoszenie incydentu
  & Pełna aktualizacja wczesnego ostrzeżenia z~wstępną oceną incydentu,
    jego dotkliwości i~skutków oraz wskaźnikami integralności systemu.
    Dostawcy usług zaufania dokonują zgłoszenia w~ciągu 24~h. \\[4pt]
Na wniosek organu: sprawozdanie okresowe
  & Aktualizacje statusu w trakcie obsługi incydentu, stosownie do żądań
    CSIRT lub właściwego organu. \\[4pt]
Max. 1 miesiąc: sprawozdanie końcowe
  & Szczegółowy opis incydentu, rodzaj zagrożenia i~pierwotna przyczyna,
    zastosowane środki ograniczające ryzyko, transgraniczne skutki incydentu.
    Jeżeli incydent trwa - w~terminie miesiąca od jego zakończenia. \\
\hline
\end{tabular}
\end{table}

Podmioty mają również obowiązek poinformowania odbiorców usług o poważnych incydentach mogących negatywnie wpłynąć na świadczone usługi, a~także
o~poważnych cyberzagrożeniach i~środkach zaradczych, jakie mogą podjąć. Samo zgłoszenie incydentu nie nakłada na podmiot zgłaszający zwiększonej
odpowiedzialności prawnej.

\subsection{Nadzór i~egzekwowanie przepisów}

Dyrektywa NIS2 wprowadza wyraźne rozróżnienie reżimów nadzoru w zależności
od kategorii podmiotu.

\begin{table}[H]
\centering
\caption{Różnice w nadzorze nad podmiotami kluczowymi i~ważnymi}
\label{tab:nis2_nadzor}
\begin{tabular}{>{\bfseries}p{4.5cm} p{4.5cm} p{4.5cm}}
\hline
Aspekt & Podmioty kluczowe & Podmioty ważne \\
\hline
Charakter nadzoru
  & Nadzór ex ante (prewencyjny, ciągły) i~ex post
  & Nadzór ex post (reaktywny, po incydencie lub sygnale) \\[4pt]
Inspekcje i~audyty
  & Regularne inspekcje na miejscu, zdalne audyty, testy
  & Inspekcje na żądanie, po stwierdzeniu naruszeń \\[4pt]
Maks. kara administracyjna
  & Do 10~mln~EUR lub 2\% rocznego światowego obrotu
  & Do 7~mln~EUR lub 1,4\% rocznego światowego obrotu \\[4pt]
Odpowiedzialność osobista
  & Możliwy zakaz pełnienia funkcji kierowniczych przez osoby fizyczne
  & Możliwy zakaz pełnienia funkcji kierowniczych przez osoby fizyczne \\
\hline
\end{tabular}
\end{table}

\subsection{Odpowiedzialność zarządcza}

Jedną z~najistotniejszych zmian wprowadzanych dyrektywą NIS2 jest \textit{bezpośrednia odpowiedzialność organów zarządzających} za przestrzeganie wymogów cyberbezpieczeństwa. Oznacza to, że członkowie zarządu i~inne osoby pełniące funkcje kierownicze mogą być osobiście pociągnięci do odpowiedzialności
za poważne naruszenia przepisów dyrektywy. Organy zarządzające podmiotów zobowiązane są do: 
\begin{enumerate}[label=(\roman*)]
  \item zatwierdzania środków zarządzania ryzykiem,
  \item nadzoru nad ich wdrożeniem,
  \item odbywania regularnych szkoleń z~zakresu cyberbezpieczeństwa. 
\end{enumerate}

W~skrajnych przypadkach właściwy organ może zwrócić się do właściwego sądu o~tymczasowe zawieszenie osoby fizycznej
w~pełnieniu funkcji kierowniczej w~danym podmiocie.


\subsection{Rejestracja i~samoidentyfikacja}

Dyrektywa NIS2 wprowadza mechanizm {samoidentyfikacji}: to sam podmiot, a nie organ nadzorczy, ma obowiązek ustalić, czy mieści się w zakresie
stosowania dyrektywy, a następnie dokonać rejestracji u właściwego organu, przekazując dane adresowe, kontaktowe, sektor i~podsektor oraz wykaz państw
członkowskich, w których świadczy usługi. O~każdej zmianie tych danych podmiot zobowiązany jest do powiadomienia organu w terminie dwóch tygodni.
W polskiej implementacji dokonanej Krajową Ustawą o Cyberbezpieczeństwie\autocite{uksc2026} termin na zgłoszenie przedsiębiorstwa do wykazu
podmiotów kluczowych i~ważnych wynosi 6~miesięcy od wejścia w~życie ustawy, zaś na wdrożenie środków zarządzania ryzykiem 12~miesięcy.

\section{Regulacje związane z rynkiem finansowym}

Rynek finansowy w Polsce podlega nadzorowi Komisji Nadzoru Finansowego (\textit{KNF}). Do stycznia 2025~roku w kontekście bezpieczeństwa i~odporności cyfrowej,
podstawowym dokumentem obowiązującym banki była {Rekomendacja D KNF}\autocite{knf_rekomendacja_d_2013}.
KNF wydaje rekomendacje na podstawie art.~137 pkt~5 ustawy Prawo bankowe z dnia 29~sierpnia 1997~roku, Rekomendacje to dokumenty o charakterze nadzorczym,
działające na zasadzie \blockquote{zastosuj albo wyjaśnij}, które efiniuje i~opisują oczekiwania KNF wobec instytucji finansowych w zakresie
zarządzania ryzykiem, bezpieczeństwa IT oraz ochrony danych klientów. Formalnie nie są aktami prawnymi, w praktyce traktowane są przez nadzorowane podmioty jako  wymogi, ponieważ ich niestosowanie skutkuje negatywną oceną w ramach Badania i~Oceny Nadzorczej (BION)\autocite{knf_bion_2024}.

Rekomendacja D została przyjęta uchwałą KNF nr~7/2013 z dnia 8~stycznia 2013~roku, następnie  uchylona z dniem 17~stycznia 2025~r. i~zastąpiona
bezpośrednio stosowanym \textit{Rozporządzeniem Parlamentu Europejskiego i~Rady (UE) 2022/2554 z dnia 14~grudnia 2022~r. w sprawie operacyjnej odporności cyfrowej sektora finansowego i~zmieniającym rozporządzenia (WE) nr~1060/2009, (UE) nr~648/2012, (UE) nr~600/2014, (UE) nr~909/2014
oraz (UE) 2016/1011}, powszechnie zwanym DORA (\textit{Digital Operational Resilience Act}) \autocite{dora2022}.

Wytyczne i~dobre praktyki zawarte w Rekomendacji~D zostały w większości zachowane i~przeniesione do DORA oraz towarzyszących jej regulacyjnych
standardów technicznych (RTS) i~wykonawczych standardów technicznych (ITS) opracowanych przez Europejski Urząd Nadzoru Bankowego (EBA).

\subsection{Prawo bankowe}
Ustawa z dnia 29 sierpnia 1997 r. - {Prawo bankowe} to akt prawny uchwalony przez Sejm RP 29 sierpnia 1997~r., opublikowana w Dzienniku Ustaw 1997 nr 140 poz. 939, weszła w życie 1 stycznia 1998~roku. Od tamtej pory była wielokrotnie nowelizowana, aktualny tekst jednolity obejmuje zmiany wprowadzone m.in. przez przepisy implementujące dyrektywy unijne (CRD IV, PSD2, AML).

Art. 104 ustawy Prawo bankowe nakłada na bank, jego pracowników oraz wszystkich pośredników wykonujących czynności bankowe bezwzględny obowiązek zachowania tajemnicy bankowej, obejmuje ona wszelkie informacje o czynnościach bankowych uzyskane podczas negocjacji, zawierania i~realizacji umowy z klientem, a więc m.in. salda rachunków, historię transakcji, dane o kredytach, zabezpieczeniach czy zdolności kredytowej. Ujawnienie tych informacji osobom trzecim jest dopuszczalne wyłącznie za pisemną zgodą klienta lub w ściśle określonych przypadkach przewidzianych przez prawo (np. na żądanie prokuratury lub organów podatkowych). W praktyce oznacza to, że każdy podmiot zewnętrzny, w tym dostawca chmury czy firma technologiczna któremu bank powierza przetwarzanie danych, przejmuje obowiązek zachowania tajemnicy bankowej i~musi gwarantować jej nienaruszalność kontraktowo oraz technicznie.

\section{Rozporządzenie DORA: cyfrowa odporność operacyjna sektora finansowego}

Rozporządzenie Parlamentu Europejskiego i~Rady (UE) 2022/2554, znane jako {DORA} (\textit{Digital Operational Resilience Act}), weszło w~życie 16~stycznia 2023~roku, a~jego stosowanie stało się obowiązkowe od 17~stycznia 2025~roku.
Rozporządzenie stanowi odpowiedź na rosnącą zależność sektora finansowego od technologii informacyjno-komunikacyjnych (ICT) oraz systematyczny wzrost liczby cyberincydentów dotykających instytucje finansowe w~całej Unii Europejskiej~\autocite{dora2022}.

Rozporządzenie strukturyzuje obowiązki wokół pięciu wzajemnie powiązanych filarów, opisanych szczegółowo poniżej.

\subsection*{Filar I, Zarządzanie ryzykiem ICT}

Pierwszy filar zobowiązuje instytucje finansowe do ustanowienia kompleksowych i~udokumentowanych ram zarządzania ryzykiem ICT, za które odpowiedzialność ponosi
bezpośrednio kierownictwo wyższego szczebla~\autocite{dora2022}. Ramy te muszą obejmować szczegółowy rejestr wszystkich aktywów ICT wspierających funkcje krytyczne, polityki bezpieczeństwa informacji, mechanizmy ciągłego monitorowania anomalii oraz plany ciągłości działania i~odtwarzania po awarii. Podmiot finansowy zobowiązany jest do regularnego przeglądu i~aktualizacji tych ram, a~organ nadzoru może w~każdej chwili
zażądać ich udostępnienia~\autocite{resilia_dora}.

\subsection*{Filar II, Zarządzanie incydentami ICT i~ich raportowanie}

Drugi filar wprowadza ujednolicone, rygorystyczne reguły klasyfikowania, rejestrowania i~zgłaszania incydentów związanych z~ICT. W~przypadku wykrycia
poważnego incydentu podmiot finansowy zobowiązany jest do przekazania wstępnego powiadomienia właściwemu organowi nadzorczemu w~ciągu 4~godzin,
raportu pośredniego w~ciągu 72~godzin oraz raportu końcowego w~ciągu miesiąca.
Naruszenie bezpieczeństwa ICT prowadzące do wycieku danych osobowych podlega równolegle obowiązkowi zgłoszenia do UODO na podstawie art.~33 RODO, co~czyni z~DORA i~RODO komplementarne reżimy sprawozdawcze.

\subsection*{Filar III, Testowanie cyfrowej odporności operacyjnej}

Trzeci filar nakłada obowiązek regularnego, ustrukturyzowanego testowania odporności systemów ICT na zagrożenia. Co najmniej raz w~roku podmioty finansowe
muszą przeprowadzać testy obejmujące m.in.\ oceny podatności, testy scenariuszowe i~testy penetracyjne. Duże, systemowo istotne instytucje zobowiązane są do wykonywania zaawansowanych testów penetracyjnych sterowanych zagrożeniami (\textit{Threat-Led Penetration Testing, TLPT}) z~udziałem zewnętrznych certyfikowanych podmiotów, przy czym zewnętrzny dostawca ICT musi w~pełni współpracować przy ich przeprowadzaniu.

\subsection*{Filar IV, Zarządzanie ryzykiem zewnętrznych dostawców ICT}

Czwarty filar stanowi najbardziej przełomowy element rozporządzenia. DORA nakłada bezpośrednie obowiązki na zewnętrznych dostawców ICT, w~tym hiperscalerów chmurowych takich jak Google Cloud, Microsoft Azure czy Amazon Web Services, obejmując ich nadzorem jako dostawców kluczowych (ang. \textit{Critical ICT Third-Party Providers
(CTPP)}). Instytucje finansowe muszą przeprowadzać ocenę ryzyka przed zawarciem każdej umowy z~dostawcą ICT, a~sama umowa musi zawierać postanowienia
o~prawie do audytu, lokalizacji przetwarzania danych, planach wyjścia (ang. \textit{exit plan}) oraz minimalnych okresach przejściowych przy zmianie dostawcy~\autocite{eiopa_dora}.
Korzystanie z~niezatwierdzonych narzędzi chmurowych (ang. \textit{shadow IT}) stanowi naruszenie tego filaru, a~jednocześnie naruszenie zasady rozliczalności z~art.~5 ust.~2 RODO.

\subsection*{Filar V, Wymiana informacji o~cyberzagrożeniach}

Piąty filar zachęca, a~w~określonych przypadkach zobowiązuje, podmioty finansowe do uczestnictwa w~schematach wymiany informacji o~zagrożeniach cybernetycznych, podatnościach i~technikach ataku w~ramach zaufanych społeczności sektorowych~\autocite{eiopa_dora}.
Celem jest budowanie zbiorowej odporności sektora: wczesne ostrzeżenie jednej instytucji o~nowym wektorze ataku pozwala innym podmiotom na podjęcie środków zapobiegawczych zanim incydent zdąży się rozprzestrzenić. Uczestnictwo w~takich schematach musi jednak odbywać się zgodnie z~przepisami o~ochronie danych osobowych oraz z~zachowaniem tajemnicy zawodowej~\autocite{resilia_dora}.

\subsection*{Sankcje}

Za naruszenie rozporządzenia Wiodący Organ Nadzorczy (ang \textit{Lead Overseer}) może nakładać na kluczowych dostawców ICT kary w~wysokości do 1\% średniego dziennego światowego obrotu za każdy dzień naruszenia, przez okres do sześciu miesięcy. Dla podmiotów finansowych nadzorowanych przez KNF sankcje określa prawo krajowe wdrażające uprawnienia nadzorcze wynikające z~rozporządzenia~\autocite{knf_dora}.

Wiodący Organ Nadzorczy to instytucja wyznaczona na poziomie UE do bezpośredniego nadzoru nad kluczowymi dostawcami ICT  w ramach rozporządzenia DORA.
Rolę Wiodącego Organu Nadzorczego pełni jedna z trzech unijnych instytucji:
\begin{enumerate}[label=(\roman*)]
  \item EBA (European Banking Authority) - dla dostawców ICT obsługujących głównie sektor bankowy,
\item ESMA (European Securities and Markets Authority) - dla dostawców obsługujących rynki kapitałowe,
\item EIOPA (European Insurance and Occupational Pensions Authority) - dla dostawców obsługujących ubezpieczenia i fundusze emerytalne.
\end{enumerate}
Przypisanie konkretnego dostawcy do jednego z tych organów zależy od tego, jaki typ podmiotów finansowych jest jego głównym klientem.
Wiodący Organ Nadzorczy ma szerokie uprawnienia wobec krytycznych dostawców ICT: może żądać informacji i dokumentacji, przeprowadzać inspekcje, wydawać zalecenia naprawcze oraz nakładać kary finansowe.

\section{Ochrona danych medycznych, regulacje prawne i~nadzór}

Dane dotyczące zdrowia należą do szczególnej kategorii danych osobowych w~rozumieniu art.~9 rozporządzenia (UE) 2016/679 (\textit{RODO})\autocite{rodo}, co oznacza, że ich przetwarzanie jest co~do zasady zakazane, a~wyjątki od tego zakazu są ściśle i~enumeratywnie określone w~przepisach~\autocite{przewodnik_rodo_zdrowie}. Do kategorii tej zalicza się wszelkie informacje o~stanie zdrowia fizycznego lub psychicznego osoby, dane genetyczne i~biometryczne, a~także informacje pośrednie (takie jak numer recepty, historia wizyt lekarskich czy przynależność do grupy ubezpieczenia zdrowotnego)  jeżeli pozwalają one na wyciągnięcie wniosków o~stanie zdrowia danej osoby.
Podwyższony reżim ochrony danych zdrowotnych uzasadniony jest ich szczególną wrażliwością: ujawnienie informacji o~chorobie, leczeniu psychiatrycznym czy uzależnieniu może prowadzić do dyskryminacji w~zatrudnieniu, wykluczenia z~rynku ubezpieczeń lub stygmatyzacji społecznej.

\subsection*{Krajowe ramy prawne i~wielość nadzorców}

Ochrona danych medycznych w~Polsce regulowana jest przez trzy nakładające się
warstwy prawa, za których egzekwowanie odpowiadają odrębne organy nadzorcze.

\begin{enumerate}[label=(\roman*)]
  \item Pierwszą warstwę stanowi RODO wraz z~krajową ustawą o~ochronie danych osobowych z~2018~roku. Nadzór nad jej przestrzeganiem sprawuje Prezes Urzędu Ochrony Danych Osobowych (UODO), uprawniony do przeprowadzania kontroli w~placówkach medycznych, wydawania decyzji nakazowych i~nakładania kar administracyjnych sięgających 20~milionów euro lub 4\% rocznego obrotu. Każda placówka medyczna przetwarzająca dane zdrowotne na dużą skalę, zobowiązana jest do wyznaczenia Inspektora Ochrony Danych (IOD), który monitoruje zgodność przetwarzania z~RODO, szkoli personel i~pełni rolę punktu kontaktowego dla pacjentów oraz organu nadzoru~\autocite{przewodnik_rodo_dane_zdrowotne}.

  \item Drugą warstwę tworzy ustawa z~dnia 6~listopada 2008~roku o~prawach pacjenta i~Rzeczniku Praw Pacjenta. Nadzór nad prawem pacjenta do dostępu do dokumentacji medycznej (art.~23--29 ustawy) sprawuje Rzecznik Praw Pacjenta, który rozpatruje skargi pacjentów i~może kierować zalecenia do podmiotów leczniczych. Ustawa nakłada na każdy podmiot udzielający świadczeń zdrowotnych obowiązek prowadzenia, przechowywania i~udostępniania dokumentacji medycznej w~ściśle określony sposób, a~dokumentacja ta podlega ochronie zarówno na podstawie ustawy, jak i~RODO~\autocite{rpp_dokumentacja}.

  \item Trzecią warstwę stanowi system akredytacji i~nadzoru nad podmiotami leczniczymi, który sprawuje {Ministerstwo Zdrowia} we współpracy z~\textit{Narodowym Funduszem Zdrowia} oraz, w~zakresie standardów informatycznych, {Centrum e-Zdrowia}, odpowiadające za systemy P1. Systemy P1 to centralny system e-zdrowia w Polsce, formalnie określany jako \textit{Elektroniczna Platforma Gromadzenia, Analizy i~Udostępniania Zasobów Cyfrowych o Zdarzeniach Medycznych}. Jest to państwowa platforma teleinformatyczna rozwijana przez Centrum e-Zdrowia, służąca do gromadzenia, przetwarzania i~udostępniania danych o zdarzeniach medycznych oraz do obsługi usług takich jak e-recepta, e-skierowanie i~wymiana EDM ~\autocite{sc_prawa_pacjenta}.
\end{enumerate}  

\subsection*{Europejska Przestrzeń Danych Zdrowotnych (EHDS)}

Nową regulacją jest \textit{Europejska Przestrzeń Danych Zdrowotnych} (\textit{European Health Data Space, EHDS}), której rozporządzenie weszło w~życie 26~marca 2025~roku~\autocite{ehds_gov_zdrowie}. EHDS tworzy
ogólnounijne ramy pierwotnego i~wtórnego wykorzystania elektronicznych danych zdrowotnych: pacjent uzyskuje prawo dostępu do swoich danych w~standardowym europejskim formacie (\textit{EEHRxF}), możliwość kontrolowania kto ma do nich wgląd oraz prawo do ograniczenia dostępu do wybranych części dokumentacji~\autocite{ehds_komisja}.
Pełne wdrożenie systemu  przewidziane jest do 26~marca 2027~roku~\autocite{ehds_gov_zdrowie}.

EHDS działa równolegle do RODO, nie zastępując go, lecz nadając sektorowi zdrowotnemu dodatkowy, sektorowy wymiar regulacyjny. Oznacza to, że placówki medyczne będą podlegać jednocześnie: RODO (nadzór UODO), ustawie o~prawach pacjenta (nadzór Rzecznika Praw Pacjenta) oraz EHDS (nadzór krajowego organu ds.\ elektronicznych danych zdrowotnych, który Polska wyznaczy do 2027~roku).
Naruszenie bezpieczeństwa systemu ICT przechowującego dane medyczne będzie tym samym incydentem, który należy raportować do wielu organów jednocześnie, analogicznie do zbieżności DORA i~RODO w~sektorze finansowym.

